\chapter{Estado da Arte}
\label{cap:estado-arte}

\section{Considerações Iniciais}

Este capítulo apresenta o estado da arte sobre a comparação entre as arquiteturas monolítica e
de microsserviços. Foram selecionados trabalhos que compararam as duas abordagens de forma
experimental ou analítica, incluindo artigos publicados em eventos e periódicos internacionais
e trabalhos de conclusão de curso desenvolvidos em instituições brasileiras. Os trabalhos são
apresentados em ordem cronológica, cada um em uma seção, destacando o objetivo, a aplicação
utilizada, os aspectos avaliados e os principais resultados. Ao final, a
Seção~\ref{sec:analise-estado-arte} compara esses trabalhos entre si e com a proposta desta
monografia.

\section{Villamizar et al. (2015)}

\citeonline{villamizar2015}, da Universidad de los Andes, na Colômbia, partiram da observação
de que grandes empresas da Internet, como Amazon, Netflix e LinkedIn, passaram a utilizar
microsserviços para implantar aplicações de grande porte na nuvem como um conjunto de pequenos
serviços, que podem ser desenvolvidos, testados, implantados, escalados e atualizados de forma
independente. O objetivo dos autores foi avaliar se esse padrão é uma boa alternativa para a
implantação de aplicações Web em nuvem.

Para isso, os autores apresentaram um estudo de caso no qual uma mesma aplicação corporativa
foi desenvolvida e implantada em nuvem de duas formas, uma monolítica e outra em
microsserviços, ambas com o \textit{framework} Play. Foram executados testes de desempenho nas
duas versões, e o trabalho discute os benefícios e os desafios que empresas já estabelecidas
podem enfrentar ao adotar microsserviços.

Em trabalho posterior, o mesmo grupo comparou o custo de infraestrutura da aplicação em três
cenários com as mesmas condições de escala: monolítica, em microsserviços operados pelo próprio
cliente e em microsserviços operados pelo provedor de nuvem. Os resultados mostraram que os
microsserviços podem reduzir o custo de infraestrutura em relação ao monólito e que o uso de
serviços de nuvem projetados especificamente para implantar e escalar microsserviços reduziu
os custos em 70\% ou mais \cite{villamizar2016}. Esses trabalhos mostram que a comparação entre
as arquiteturas não se limita ao desempenho, envolvendo também o custo e a forma de
implantação.

\section{Al-Debagy e Martinek (2018)}

\citeonline{aldebagy2018}, da Universidade de Tecnologia e Economia de Budapeste, motivaram seu
estudo pelas contradições encontradas na literatura: enquanto alguns trabalhos relatavam maior
vazão para os microsserviços, outros apontavam desempenho superior do monólito em diversas
configurações. Para investigar essas diferenças, os autores utilizaram a plataforma JHipster
para gerar aplicações com Spring Boot e AngularJS nas duas arquiteturas. A versão em
microsserviços era composta por três serviços, um registro de serviços e um
\textit{gateway}, e todas as aplicações foram executadas em contêineres Docker, em um servidor
com 16~GB de memória, enquanto a ferramenta Apache JMeter gerava a carga a partir de um cliente
remoto.

Foram definidos três cenários de teste, avaliados pela vazão e pelo tempo de resposta. No teste
de carga, o número de usuários simultâneos foi aumentado gradualmente de 100 até 7.000: as duas
arquiteturas tiveram desempenho semelhante, com diferença média de 0,87\%, embora o monólito
tenha se saído melhor com poucos usuários. No teste de concorrência, em que todos os serviços
eram acionados ao mesmo tempo com 100 a 1.000 requisições cada, o monólito apresentou vazão
cerca de 6\% superior, sem diferença significativa no tempo de resposta. No terceiro cenário,
com 10.000 usuários simultâneos, os autores compararam duas tecnologias de descoberta de
serviços na versão em microsserviços: a configuração com Consul obteve vazão cerca de 3,8\%
maior e tempo de resposta ligeiramente menor do que a configuração com Eureka.

Os autores concluíram que as duas arquiteturas têm desempenho semelhante sob carga normal e
recomendaram o monólito para aplicações pequenas, utilizadas por poucos usuários. O trabalho
também evidencia que as escolhas de infraestrutura da arquitetura distribuída, como a
ferramenta de descoberta de serviços, influenciam o resultado da comparação.

\section{Mendes (2021)}

\citeonline{mendes2021}, em trabalho de conclusão do curso de Engenharia de Software da
Universidade de Brasília, partiu do contexto de empresas que migram seus sistemas para
microsserviços em busca dos benefícios dessa arquitetura, mas que, quando o fazem sem
planejamento adequado e em desacordo com os objetivos de negócio, podem levar o novo modelo
arquitetural ao fracasso.

O trabalho consistiu em uma pesquisa exploratória sobre os estilos monolítico e de
microsserviços, a partir da qual foi construído um mapa mental com os fatores que afetam cada
modelo arquitetural. Esses fatores foram então analisados em quatro estudos de caso de
empresas que migraram de um modelo arquitetural para o outro.

A autora concluiu que a arquitetura monolítica é indicada para a fase de descoberta do domínio,
mas tende a perder manutenibilidade e evolucionabilidade à medida que a base de código cresce,
enquanto a arquitetura de microsserviços tende a ser mais sustentável, porém apresenta custo
elevado e exige domínio sobre o problema e as tecnologias envolvidas. Trata-se de uma análise
qualitativa, sem a implementação e a medição de aplicações.

\section{Blinowski, Ojdowska e Przybyłek (2022)}

\citeonline{blinowski2022} observaram que empresas menores costumam esperar, ao migrar para
microsserviços, os mesmos ganhos obtidos por grandes empresas, o que pode não se confirmar,
especialmente em sistemas sem altíssima concorrência que podem ser escalados verticalmente.
Como as evidências sobre o tema eram escassas e inconsistentes, os autores se propuseram a
comparar o desempenho e a escalabilidade das duas arquiteturas por meio de uma aplicação Web de
referência.

A aplicação foi implementada em quatro versões, combinando dois estilos arquiteturais
(monolítico e de microsserviços) e duas tecnologias (Java e C\#.NET). Foram realizados
experimentos controlados em três ambientes de implantação: uma máquina local e os serviços
Azure Spring Cloud e Azure App Service, da plataforma de nuvem Microsoft Azure.

Os resultados indicaram que: em uma única máquina, o monólito teve desempenho superior à
versão em microsserviços; a tecnologia Java aproveitou melhor máquinas mais potentes em
serviços com processamento intensivo, enquanto C\#.NET se saiu melhor em serviços leves
executados em máquinas de baixa capacidade; a escalabilidade vertical foi mais vantajosa em
custo do que a horizontal na nuvem Azure; acrescentar instâncias além de certo ponto passou a
degradar o desempenho; e a tecnologia de implementação não teve impacto perceptível na
escalabilidade. O trabalho reforça que o ambiente de execução e a estratégia de escala
influenciam diretamente a comparação entre as arquiteturas.

\section{Santos Junior (2024)}

\citeonline{santosjunior2024}, em trabalho de especialização em Engenharia de Software da
Pontifícia Universidade Católica de São Paulo, relatou a experiência de um projeto integrador
voltado à migração de uma aplicação de \textit{e-commerce} monolítica para microsserviços.

A parte teórica do trabalho compara as arquiteturas monolítica e de microsserviços e discute os
desafios da migração entre elas. A parte prática descreve a arquitetura adotada para atender às
regras de negócio do \textit{e-commerce}, incluindo uma camada de \textit{gateway}, com o
objetivo de obter maior escalabilidade, resiliência e manutenibilidade. O código-fonte do
projeto foi disponibilizado publicamente.

O trabalho se aproxima desta monografia por tratar do mesmo domínio de aplicação, mas tem
caráter de relato de experiência: a proposta de migração não é acompanhada de uma comparação
quantitativa entre as duas versões.

\section{Vale (2025)}

\citeonline{vale2025}, em trabalho de conclusão do curso de Engenharia de Computação da
Universidade Federal do Ceará, investigou o debate entre as arquiteturas monolítica e de
microsserviços buscando passar da análise teórica para a validação empírica. Para isso,
construiu dois protótipos funcionalmente idênticos de um sistema de cadastro, um monolítico e
outro em microsserviços, desenvolvidos em JavaScript com Node.js e Express e executados em
contêineres Docker.

Os protótipos foram comparados quanto ao tempo de resposta, ao consumo de processador e de
memória e à complexidade de implantação. O monólito apresentou menor latência e infraestrutura
mais simples, enquanto a versão em microsserviços registrou latência cerca de 147\% maior em
uma operação que envolvia vários serviços, consumo de memória em repouso cerca de 271\% maior
e mais que o triplo de artefatos de configuração.

Entre os trabalhos analisados, este é o que mais se aproxima da metodologia proposta nesta
monografia, por comparar duas versões funcionalmente equivalentes de uma mesma aplicação em
ambiente controlado. A aplicação avaliada, entretanto, é um sistema de cadastro com poucas
funcionalidades, o que limita a observação de aspectos como comunicação assíncrona,
integração com serviços externos e consistência de dados entre serviços.

\section{Análise Comparativa dos Trabalhos}
\label{sec:analise-estado-arte}

A Tabela~\ref{tab:trabalhos-relacionados} resume os trabalhos apresentados neste capítulo e
posiciona a presente monografia em relação a eles.

\begin{table}[htbp]
\centering
\caption{Resumo dos trabalhos do estado da arte}
\label{tab:trabalhos-relacionados}
\small
\begin{tabular}{|>{\raggedright\arraybackslash}p{2.9cm}|>{\raggedright\arraybackslash}p{3.4cm}|>{\raggedright\arraybackslash}p{3.6cm}|>{\raggedright\arraybackslash}p{5.0cm}|}
\hline
\textbf{Trabalho} & \textbf{Aplicação e tecnologia} & \textbf{Aspectos avaliados} & \textbf{Principais resultados} \\ \hline
\citeonline{villamizar2015} & Aplicação corporativa, Play \textit{Framework}, nuvem & Desempenho, implantação e custo em nuvem & Benefícios e desafios da adoção de microsserviços; redução de custo de infraestrutura \\ \hline
\citeonline{aldebagy2018} & Aplicações geradas com JHipster (Spring Boot), Docker & Vazão e tempo de resposta sob carga e concorrência & Desempenho semelhante sob carga; monólito com vazão cerca de 6\% maior em concorrência \\ \hline
\citeonline{mendes2021} & Pesquisa exploratória e quatro estudos de caso & Manutenibilidade e evolucionabilidade & Monólito indicado para descoberta do domínio; microsserviços mais sustentáveis, porém mais custosos \\ \hline
\citeonline{blinowski2022} & Aplicação Web em Java e C\#.NET, local e Azure & Desempenho e escalabilidade & Monólito mais rápido em uma máquina; escala vertical mais vantajosa; excesso de instâncias degrada o desempenho \\ \hline
\citeonline{santosjunior2024} & \textit{E-commerce}, migração com \textit{gateway} & Escalabilidade, resiliência e manutenibilidade & Proposta de arquitetura para a migração, sem medição comparativa \\ \hline
\citeonline{vale2025} & Protótipo de cadastro, Node.js e Docker & Tempo de resposta, CPU, memória e implantação & Monólito com menor latência e infraestrutura mais simples \\ \hline
Este trabalho & \textit{E-commerce} com pagamentos e mensageria, Java e Spring Boot & Desempenho, escalabilidade, manutenibilidade, complexidade e implantação & Monólito mais rápido, mais econômico e mais simples, e com escala horizontal eficiente; microsserviços com isolamento de falhas e escala por serviço \\ \hline
\end{tabular}
\par\vspace{4pt}{\footnotesize Fonte: Elaborada pelo autor (2026).}
\end{table}

A análise dos trabalhos mostra que os resultados não são uniformes. Em geral, os estudos
experimentais apontam vantagem do monólito em desempenho quando executado em uma única máquina
ou sob carga moderada \cite{aldebagy2018, blinowski2022, vale2025}, enquanto os benefícios dos
microsserviços aparecem em aspectos como custo de infraestrutura em nuvem, implantação
independente e sustentabilidade da evolução do código \cite{villamizar2016, mendes2021}. Os
resultados dependem, portanto, da aplicação, da carga, do ambiente de execução e das escolhas
de infraestrutura, o que reforça a importância de comparações controladas.

Observa-se também que os trabalhos experimentais utilizaram aplicações simplificadas, como
sistemas de cadastro ou aplicações geradas automaticamente, e avaliaram principalmente o
desempenho, enquanto os trabalhos que trataram de manutenibilidade e evolução o fizeram de
forma qualitativa. O presente trabalho se diferencia por comparar as duas arquiteturas em uma
aplicação de \textit{e-commerce} com funcionalidades realistas, organizada em seis módulos e
incluindo autenticação com JWT, pagamentos integrados ao Mercado Pago em um módulo próprio,
armazenamento de imagens e comunicação assíncrona com Apache Kafka, e por garantir, por meio
de testes automatizados, que as duas versões apresentam comportamento funcional idêntico. Além
disso, a avaliação abrange simultaneamente os cinco atributos de qualidade descritos no
Capítulo~\ref{cap:referencial}, combinando medições de desempenho e de escalabilidade com
métricas de código e de implantação.

\section{Considerações Finais}

Este capítulo apresentou o estado da arte sobre a comparação entre as arquiteturas monolítica
e de microsserviços. Os trabalhos analisados mostram que não há consenso sobre qual
arquitetura é superior: o monólito tende a apresentar melhor desempenho em cenários de menor
escala, enquanto os microsserviços oferecem vantagens de custo, implantação e evolução em
contextos específicos. Essa divergência, somada ao uso de aplicações simplificadas e à
avaliação de poucos atributos por estudo, justifica a comparação proposta neste trabalho, cuja
metodologia é descrita no Capítulo~\ref{cap:metodologia}.
